\section{Introduction \& scope}\label{sec:intro}

QtRocket is an open-source model-rocket simulator built as a shared C++23 simulation core
with three front ends: a Qt6 Widgets GUI (\code{qtrocket}), a Qt-free headless REPL
(\code{qtrocket-cli}), and a standalone OpenGL design viewer (\code{qtrocket-visualizer}).
This document is a whole-system architecture review: a survey of what the program can
actually do today, how its layers fit together, and --- measured against OpenRocket as a
reference point --- what it cannot do yet. It is written for maintainers and future
contributors who need an honest map of the territory before choosing what to build next.

Two framing points define the document's stance. First, QtRocket is an \emph{independent}
project with design goals similar to OpenRocket's, not a clone of it; OpenRocket is used
throughout as a comparison baseline because it is the de facto standard for what a mature
open-source model-rocket simulator offers, not because feature parity is the goal
(\cref{sec:comparison}). Second, the review reports the \emph{code}, not the project's
aspirations: design documents under \srcf{docs/} were treated as claims to test against
the source, never as facts to inherit, and where the two disagree this document says so.

\subsection{The tree under review}

Every statement in this document is pinned to one revision: commit \code{254cd41}
(``Major comment cleanup''), reviewed with a clean working tree on 2026-07-06. That
commit is the tip of the \code{development} line at the time of review and includes the
completed Part Placement feature (the station-pair --- \code{StationLink} --- placement backbone specified in the
companion whitepaper under \srcf{docs/PartPlacementDesignLatex/}). Two commits of
in-flight GUI work exist beyond it on the \code{RocketTreeViewImplementation} branch ---
an initial part-tree widget and work-in-progress design save/load wiring --- and are
covered explicitly as future work (\cref{sec:interfaces,sec:roadmap}) rather than folded
silently into the review. Code citations take the form \src{model/parts/Part.h}{42};
line numbers will drift as the tree moves and should be read as anchors into that
revision, not absolutes.

\subsection{What the review found, in one paragraph}

The load-bearing finding is a pattern, not a single fact. QtRocket today has a genuinely
finished, well-guarded 3-DOF point-mass flight simulator --- composite part tree with
time-aware mass, CG, and inertia; validated US Standard Atmosphere; adaptive RKF45
integration; typed termination diagnostics; motor ingest from three sources; design
persistence; a scriptable REPL that exercises all of it --- and, surrounding that core, a
ring of machinery that is built, unit-tested, and consumed by nothing. The Barrowman
center-of-pressure composition, the per-step composite inertia tensor (recorded, never integrated), the wind model,
the quaternion slots in the state vector, and the \code{getTorques()} interface all
exist and all have zero production consumers; the test suite says so itself
(``nothing in production consumes this in P2,'' \src{model/tests/AeroTests.cpp}{18}).
The project's engineering discipline --- fail-closed diagnostics, strict warnings,
five-platform CI, pinned dependencies --- is well ahead of its feature completeness. The
architecture is not the obstacle to a realistic flight simulation; \emph{consumption} of
what the architecture already provides is.

\begin{keyinsight}[How to read the gap findings]
Throughout this document, ``dormant'' means built and tested but unreachable from any
production code path; ``partial'' means reachable but with named gaps; ``absent'' means
no trace in the tree (each absence claim was grounded in a failed search of the source,
not an impression). The distinctions matter for planning: consuming a dormant seam is
weeks of work where building an absent subsystem is months.
\end{keyinsight}

\subsection{Method}

The review was produced by independent parallel deep-reads of eight subsystem areas ---
the part tree and placement system, the motor subsystem, the simulation core,
aerodynamics and flight fidelity, the GUI and visualizer, the CLI and persistence, build
and test infrastructure, and a dedicated cross-cutting incompleteness sweep --- each
grounded in a full read of its files and each followed by an adversarial verification
pass that re-checked every file:line citation, every quoted snippet, and every
implemented/dormant/absent status claim against the source. Two behaviors were
confirmed by execution rather than reading: the CLI command inventory and an end-to-end
flight through the built REPL binary. The OpenRocket baseline of \cref{sec:comparison}
was researched from primary sources and reconciled against the verified findings.
\Cref{sec:repro} records the full method and the document's own build pipeline.

\subsection{Document roadmap}

\Cref{sec:overview} maps the system: the three executables, the strict layering, and the
build that enforces it. \Cref{sec:model,sec:motors,sec:sim} review the three pillars of
the core --- the part tree with its placement resolver and mass-property machinery, the
motor subsystem from thrust curves to the three ingest paths, and the propagator with its
integrators and environment models. \Cref{sec:aero} is the keystone: what the force path
actually computes, what sits dormant beside it, and what that means for flight fidelity.
\Cref{sec:interfaces} reviews the three user-facing surfaces and their very different
levels of completeness; \cref{sec:persistence} the design-file format and
interoperability; \cref{sec:testing} the test suites, coverage tooling, and CI.
\Cref{sec:comparison} sets the feature matrix against OpenRocket and ranks the
functional deficiencies; \cref{sec:roadmap} consolidates every incomplete area found by
the review into a severity-ordered inventory with a suggested sequencing. \Cref{sec:repro}
closes with reproducibility: how this document builds and how the review was performed.
